\begin{textsample}{NEG Dim 3 – pi – Score -2.00 – pi/pi\_em\_61.txt}  \label{ex:f3_neg_028}
Matriz Curricular A Matriz Curricular da área de Ciências Humanas e Sociais Aplicadas busca servir como organizador e um mapa de aprendizagens visando ampliar conhecimentos, através da superação das lacunas de aprendizagem que marcam o percurso formativo dos estudantes da educação básica, de modo que auxilie o professor a criar estratégias de ensino que valorize as experiências cognitiva dos estudantes.
Espera -se que o conhecimento seja significativo, a partir da integração dos conhecimentos intelectuais e práticos, com base na proposta da LDB nº 9.394/96 de desenvolver um processo educativo com base nas reais necessidades do estudante.
Sugestão reforçada pela BNCC, ao focar a centralidade do estudante e seu protagonismo na construção do conhecimento e sua atuação diante dos desafios que marcam a contemporaneidade.
Com vista a atender a progressão de aprendizagens a presente matriz curricular está organizada conforme a BNCC ( 2018 ), considerando as \textbf{competências} gerais da educação básica, associadas às \textbf{competências} específicas da área, definidas com base nos processos cognitivos e socioemocionais a serem trabalhados, as habilidades a elas inter-relacionadas para o desenvolvimento das categorias centrais da área, indicação de objetos do conhecimento e os objetivos de aprendizagem a serem desenvolvidos em cada série do Ensino Médio, de forma a facilitar as ações educacionais.
Nesse sentido, a Matriz especifica os objetos do conhecimento e objetivos de aprendizagem por componentes, para que cada um seja identificado nesse processo, mas ressalta a proposta de que estes sejam trabalhados de forma articulada e interdisciplinar, com sobreposição do trabalho por área, com vista a formação integrada.
Assim, a organização por área do conhecimento, conforme cita a BNCC ao referenciar o Parecer CNE/CP nº 11/2009²⁵ : > “ não exclui necessariamente as disciplinas, com suas especificidades e saberes próprios historicamente construídos, mas, sim, implica o fortalecimento das relações entre elas e a sua contextualização para apreensão e intervenção na realidade, requerendo trabalho conjugado e cooperativo dos seus professores no planejamento e na execução dos planos de ensino ” ( BRASIL, 2009 ; ênfases adicionadas apud BNCC, 2018, p. 32 ).
Assim, apesar dos objetos do conhecimento estarem separados por componentes, enfatiza -se que eles podem ser trabalhados por todos os componentes com abordagem específica de forma interdisciplinar, considerando o tema geral norteado por cada \textbf{competência} específica.
Partindo deste ponto de vista, os limites entre cada componente curricular acabam se mesclando como propõe a BNCC, num ensino efetivamente por área de conhecimento, requerendo trabalho conjugado e cooperativo dos seus professores no planejamento dentro da área de forma interdisciplinar e transdisciplinar.
Assim, os temas gerais permitem aos professores vislumbrarem melhor as categorias conceituais da área a serem aprofundadas, a abordagem dos temas transversais contemporâneos e o diálogo interdisciplinar com as demais áreas de conhecimento.
Ressalta -se que em decorrência da carga horária anual de 80 horas para área de Ciência Humanas e sociais aplicadas na 3ª série, somente os componentes de História e Geografia foram considerados.
Contudo, os saberes de Filosofia e Sociologia estão assegurados, em conformidade com a BNCC ( 2018 ) que faz referência aos estudos e práticas destes referidos componentes, e em acordo com a Resolução nº 3/2018 – MEC/CNE/CEB, que em seu artigo 11, inciso VIII, § 5º destaca a possibilidade dos estudos e práticas de Sociologia e Filosofia serem desenvolvidos por projetos, oficinas, laboratórios, dentre outras estratégias de ensino e aprendizagens que rompam com o trabalho isolado por disciplina.
Nessa perspectiva, Filosofia e Sociologia estão como componentes na 1ª e 2ª séries e presentes em forma de estudos e práticas na 3ª série.
De forma que a carga horária anual dos componentes da área, na Formação Geral Básica, estão assim distribuídas :

% matched lemmas: competência
\end{textsample}
